\section{Experiments}
\label{sec:experiments}

We answer nine questions. \textbf{RQ1 (safety)}: which optimizers actually
keep their contract --- verified directly on the \emph{exact risk over a
fixed finite population} across repeated calibration draws?
\textbf{RQ2 (efficiency)}: at what cost, relative to fixed plans and to
hindsight oracles --- including strict low-$\alpha$ contracts with
certified deferral? \textbf{RQ3 (calibration)}: how tightly does realized
violation track the contract level $\alpha$? \textbf{RQ4 (rewrites)}: do
risk vouchers achieve their nominal coverage over repeated draws, and how
loose is voucher composition versus direct certification?
\textbf{RQ5 (groups)}: what does marginal
certification hide on query slices, and does group-conditional
certification fix it? \textbf{RQ6 (drift)}: are workload and index-staleness
drifts detected, at full and subsampled audit rates, with what lifetime
guarantee across repeated alarms --- and what does auditing cost?
\textbf{RQ7 (multi-constraint)}: do joint four-constraint contracts
certify, and what does the intersection--union certification buy in
power over Bonferroni?
\textbf{RQ8 (generalization)}: do the guarantees and the savings survive a
change of model family, model developer, serving backend, and cost metric
--- including a local open-weight family metered in GPU-seconds?
\textbf{RQ9 (planning \& scale)}: does cardinality-driven access-path
selection beat fixed join orders, and how does certification scale to
$10^5$-candidate plan spaces?

\subsection{Setup}
\label{sec:exp-setup}

\textbf{Tracks.} (A) \textbf{HybridQA}~\cite{chen2020hybridqa}: multi-hop
QA over Wikipedia tables joined with entity passages; sources are a
structured row-selection operator and sparse/dense passage indexes over
440K passages; quality loss $\mathbf{1}[\text{token-F1} < 0.5]$; splits
train/cal/test = 5000/3000/3466 (train and cal drawn from the official
62K-query train set; test is the \emph{entire} official HybridQA-dev
release).
(B) \textbf{CRAG}~\cite{yang2024cragbench}: web QA with the official mock
knowledge-graph API (movie\slash music\slash sports\slash
finance\slash open domains, dynamism
labels static\slash slow\slash fast\slash real-time); evidence chunks
carry crawl ages;
answers are graded correct\slash missing\slash incorrect by a qwen-max
judge with the
gold answer (CRAG's protocol); loss is either \emph{correctness},
$\mathbf{1}[\text{not correct}]$, or \emph{hallucination},
$\mathbf{1}[\text{answered and wrong}]$; splits 700/1000/1006 covering
the \emph{full} 2706-query public Task~1\&2 development pool.
(C) \textbf{ASQA}~\cite{stelmakh2022asqa} under the official ALCE
evaluation~\cite{gao2023alce}: long-form answers with citations over a
pooled GTR corpus (89K passages), citation recall/precision scored by the
official TRUE-NLI T5-XXL entailment model~\cite{honovich2022true}; loss
$\mathbf{1}[\text{citation-recall} < 50]$; splits 150/350/448 covering
the \emph{full} 948-query ALCE-ASQA release: test is the complete
original held-out set, and the RQ1 protocol below additionally resamples
over the pooled 798 cal+test queries, so no safety conclusion rests on a
single small draw.
(D) \textbf{QAMPARI}~\cite{amouyal2022qampari} under the same ALCE
protocol: list answers (each entity must cite exactly one document) over
a pooled 86K-passage GTR corpus, citations scored by TRUE-NLI with the
official comma-split protocol; loss
$\mathbf{1}[\text{citation-recall} < 50]$; splits 150/350/500 covering
the full 1000-query ALCE-QAMPARI release. QAMPARI is the adversarial
regime for citation contracts --- even the strongest plan violates on
61\% of queries --- so it probes certification exactly where
point estimates are noisiest.

\textbf{Why these scales.} A $(C,\delta)$-certificate is a finite-sample
statement: its radius shrinks as $O(\sqrt{\log(1/\delta)/n_{\mathrm{cal}}})$.
At $n_{\mathrm{cal}}\!\in\!\{350,1000,3000\}$ and $\delta{=}0.1$ the
Hoeffding--Bentkus half-width at mid-range $\alpha$ is roughly
$0.075$/$0.04$/$0.025$ --- tight enough that certificates discriminate
among ladder rungs rather than collapsing to the strongest plan, which is
exactly the operating regime of a deployed RAG service that recalibrates
weekly on a few hundred to a few thousand audited queries. HybridQA's
test set is the full public\_dev benchmark; CRAG, ASQA, and QAMPARI
consume every publicly released query for the track (no further held-out
pool exists).
For ASQA in particular --- the smallest track --- no single-split number
is load-bearing: the RQ1 protocol resamples 1000 calibration/test
partitions from the pooled 798 queries, the reported violation
probabilities are means over those draws, and at $n_{\mathrm{cal}}=350$
the certificate's finite-sample radius ($\approx$0.075) is exactly what
the theory prescribes paying at this audit budget --- the guarantee holds
at every $n$, only its tightness scales.
Train sizes are chosen to saturate the logistic sufficiency scorers
(diminishing returns beyond a few thousand labeled executions in our
ablations). Every number below is measured on these real executions, not
on subsampled proxies.

\textbf{Ladders.} Each track has a four-rung cost-ascending ladder
generated by the rewrites of \S\ref{sec:rewrites}: from sparse-only
retrieval with qwen-flash up to hybrid RRF retrieval, cross-encoder
reranking, NLI verification, KG-API evidence (CRAG), citation repair
(ASQA, QAMPARI), and qwen-max. All executions are real API calls (28--48-way
parallel) with exact token metering; local models (BGE embeddings,
bge-reranker-v2-m3, DeBERTa-NLI, TRUE-NLI) run on two RTX 4090 GPUs.
Reported cost is LLM spend in milli-CNY per query; a query stopping at rung
$j$ is charged rungs $0..j$ cumulatively for escalation policies and only
rung $j$ for jump policies. Local GPU compute and KG-API planning calls
are metered separately and shift no conclusion
(\S\ref{sec:system}, total cost of ownership). RQ8 additionally swaps the
generator column of the HybridQA ladder across four model families from
four \emph{model developers} --- Alibaba's Qwen (flash/plus/max),
DeepSeek (v4-flash/v3.2/v4-pro), and Zhipu's GLM (4.7/5/5.2), plus
Google's open-weight Gemma~3 (1b/4b/12b/27b) --- served through two
\emph{serving backends}: the commercial families are hosted on one API
platform (DashScope) at each developer's list prices, while Gemma~3 is
served \emph{locally} by Ollama and metered in measured GPU-seconds
priced at the same RTX-4090 rental rate --- over shared materialized
retrieval (\S\ref{sec:planner}). We keep model developer and serving
backend distinct throughout: the experiment varies the family and its
price structure, not the number of API platforms.

\textbf{Baselines.} \emph{Fixed plans}: each rung as a static pipeline
(these are also the plans a workload-level optimizer chooses among).
\emph{Complexity router}: Adaptive-RAG~\cite{jeong2024adaptiverag} adapted
to our ladders --- a text classifier trained to predict the cheapest
correct rung. \emph{Utility router}: per-query
$\arg\max_j \hat{s}_j - \lambda \bar c_j$ with $\lambda$ tuned to the
constraint on tuning data (cost-aware router
style~\cite{ding2024hybridllm,hu2024routerbench}).
\emph{BO-tuned}: Optuna~\cite{akiba2019optuna} TPE over escalation
thresholds, minimizing cost s.t.\ empirical risk $\le \alpha$ on tuning
data. \emph{Abacus}~\cite{russo2025abacus}: workload-level constrained
selection with UCB sampling on a budget of 150 plan executions, choosing
the cheapest rung whose estimated risk clears $\alpha$ (point estimate).
\emph{Empirical-no-cert}: our own
ladder with the cheapest thresholds whose \emph{empirical} calibration risk
is $\le\alpha$ --- i.e., \system\ minus the statistical test (ablation).
Two hindsight skylines with different targets: the \emph{zero-loss
oracle} picks the per-query cheapest rung with no loss (a stricter
objective than any contract, hence not cost-comparable to
contract-feasible methods), and the \emph{matched oracle} is the
cheapest per-query rung assignment whose test-set mean loss is
$\le \alpha$ --- the hindsight optimum of the same constrained problem
\system\ solves, and its fair cost lower bound.
All methods consume identical execution matrices, so comparisons are
exact and reproducible.

\subsection{RQ1: contract safety across calibration draws}
\label{sec:exp-safety}

The theory bounds the probability that the deployed policy's \emph{true}
risk exceeds $\alpha$; we verify that claim directly rather than through
the binomial noise of a held-out split. For each track we treat the full
pooled cal+test execution matrix as a fixed finite population and repeat
1000 times: draw a random calibration subset, let every method
select/tune on it (scorers and routers use the fixed train split), and
evaluate the selected policy's \emph{exact} risk on the entire
population. Table~\ref{tab:safety} reports
$\Pr(\text{population risk} > \alpha)$ --- the exact finite-population
analogue of the $\delta$-guarantee ($\delta = 0.1$ throughout); at 1000
draws the Monte Carlo 95\% interval on a reported rate of $0.05$ is
$\pm 0.014$, and there is no test-sampling error in the verdict.
Calibration sets are drawn \emph{without replacement}, while the HB
p-value assumes i.i.d.\ draws; Appendix~\ref{app:finitepop} closes this
gap formally --- the hypergeometric test is exact for the protocol, the
shipped p-value is dominated by it at every configuration used here
(valid, merely conservative), and re-certifying every draw with the
exact finite-population p-value leaves all conclusions unchanged while
certifying cheaper policies.
Throughout the paper we keep three quantities distinct:
\emph{certification} is an event on the calibration draw; it purchases a
guarantee on the \emph{population} risk with failure probability
$\le \delta$ over draws; and any single test-split \emph{empirical} risk
is a noisy estimate that may sit on either side of $\alpha$ without any
guarantee being violated.

\begin{table}[t]
\caption{RQ1: probability (over 1000 calibration draws) that the deployed
policy's exact finite-population risk exceeds $\alpha$ --- the quantity
the certificate bounds --- and mean cost relative to always running the
strongest plan. Cells with no violating draw report the one-sided
95\% binomial upper bound $<.003$ rather than an absolute rate;
since each policy's population risk is one exact number,
draws vary only through selection, and $<.003$ means no draw
selected an out-of-contract policy. \system\ is the only optimizer
meeting the $\delta=0.1$
budget on all tracks.}
\label{tab:safety}
\centering\scriptsize
\setlength{\tabcolsep}{2.6pt}
\begin{tabular}{lcccccccc}
\toprule
& \multicolumn{2}{c}{HybridQA ($\alpha{=}.35$)} &
\multicolumn{2}{c}{CRAG ($\alpha{=}.65$)} &
\multicolumn{2}{c}{ASQA ($\alpha{=}.25$)} &
\multicolumn{2}{c}{QAMPARI ($\alpha{=}.66$)} \\
Method & viol. & cost & viol. & cost & viol. & cost & viol. & cost \\
\midrule
Empirical (no cert.) & .470 & .19 & .493 & .15 & .408 & .45 & .451 & .45 \\
BO-tuned             & .572 & .12 & .005 & .08 & .410 & .12 & .532 & .21 \\
Utility router       & .160 & .07 & $<$.003 & .03 & .106 & .08 & .497 & .40 \\
Abacus               & .099 & .34 & .275 & .16 & $<$.003 & .25 & .433 & .51 \\
\system\ (ladder)    & .006 & .31 & .005 & .38 & .002 & .51 & .004 & .75 \\
\system\ (full)      & \textbf{$<$.003} & .11 & \textbf{$<$.003} & .07 & \textbf{$<$.003} & .43 & \textbf{.001} & .84 \\
\bottomrule
\end{tabular}
\end{table}

Three observations. (i) Empirically tuned methods select policies whose
\emph{true} risk is out of contract in up to 49\% (empirical tuning) and
57\% (BO) of draws at these headline levels
--- genuine failures of the deployed policy, not artifacts of a
noisy test split; sweeping $\alpha$ (Fig.~\ref{fig:sweep}, thirteen
settings) shows no baseline staying below $\delta$ everywhere, with BO
reaching 64\%. The violations are not rare tail events but the typical
behavior of tuning against a noisy constraint. (ii) The certified methods
hold the population promise with margin: the full optimizer's true risk
exceeds $\alpha$ in at most 0.1\% of draws at the headline levels of
Table~\ref{tab:safety} (below $.003$ on three tracks, a single
violating draw on QAMPARI) and
never above 0.8\% anywhere on the
thirteen-setting sweep --- an order of magnitude inside the
$\delta=10\%$ budget. Rates this small are the predicted outcome, not
an artifact: certification requires clearing $\alpha$ by the
finite-sample radius (0.02--0.07 at these $n$), so an out-of-contract
deployment needs a false certification whose realized probability sits
far below its $\delta$ bound; where the 1000 draws contain no such
event we report the binomial upper bound $<.003$ --- the protocol's
resolution --- rather than an absolute rate; the escalation-only
ladder stays at or below 0.8\%. The same quantization produces the occasional near-floor
baseline entries
(utility router at CRAG's loose $\alpha = .65$: below resolution here
but 22.1\%
on the held-out criterion and far above $\delta$ at other levels in
Fig.~\ref{fig:sweep}) --- a coincidence of one setting, not a
guarantee. (On the conventional held-out criterion the certified
policies show
0--7\% rates; the difference is pure test-split binomial noise.)
(iii) Safety
is not purchased with money: on CRAG the full optimizer is \emph{both} the
safest and among the cheapest (cost ratio 0.07 vs.\ 0.15 for uncertified
empirical tuning), because the richer certified family (jump routers)
dominates threshold tuning. Abacus is safe exactly when inter-plan gaps
are large relative to its sampling noise (ASQA) and unsafe when they are
not (CRAG, 28\%; QAMPARI, 43\%) --- point estimates cannot tell these
regimes apart, which is the core argument for certification.

\smallskip\noindent\textit{Takeaway: verified on exact population risk,
point-estimate feasibility is a coin flip; certification makes the
$\delta$-promise hold with an order of magnitude to spare at every
contract level we tested, and on some tracks for free.}

\subsection{RQ2: cost at matched contracts}
\label{sec:exp-cost}

Table~\ref{tab:main} and Fig.~\ref{fig:tradeoff} (full grids in the
artifact) fix representative contract levels and show one draw each; the
guarantee behind the bold rows is the population-level one verified in
Table~\ref{tab:safety}. At $\alpha=0.34$ on HybridQA \system\ certifies a
jump-router policy costing 1.05 mCNY/query with test risk 0.338 ---
6.9$\times$ cheaper than the strongest plan (7.21) whose risk is 0.290,
and 2.0$\times$ cheaper than the certified escalation-only ladder; the
uncertified methods that undercut it (utility router at 0.40) do so by
\emph{breaking the contract} (risk 0.387~$>\alpha$). At
$\alpha=0.65$ on CRAG the certified policy costs 0.30 vs.\ 6.00 mCNY
(20$\times$); at $\alpha=0.71$ on QAMPARI, 0.43 vs.\ 6.49 mCNY
(15$\times$) --- here certification adopts the utility router's own
policy, showing that the certified family subsumes the baselines'
policies and differs only in the feasibility test; at $\alpha=0.216$ on
ASQA, 6.56 vs.\ 9.01 mCNY with test risk 0.145 while the utility router
(0.223) and BO (0.221) sit above $\alpha$. Several single-draw entries
of certified ladder policies sit slightly above $\alpha$ (0.358 vs.\
0.34 on HybridQA; 0.724 vs.\ 0.71 on QAMPARI): these are within one
binomial standard error of the test split, and Table~\ref{tab:safety}
shows the corresponding \emph{population} risks violate in at most
0.6\% of draws --- the certificate promises the distributional
guarantee, not that every finite sample lands below $\alpha$. Abacus
lands inside the contract on
ASQA at 1.31 mCNY --- but with no certificate: the identical procedure
deploys truly out-of-contract policies in 28\% of draws on CRAG
(Table~\ref{tab:safety}), and at
$n_{\mathrm{cal}}=350$ the certified margin genuinely costs the
difference; with more audit data the certified policy converges to the
same cheap plan (Proposition~\ref{prop:consistency}). The fair skyline
for \system's objective is the \emph{matched oracle} --- the cheapest
hindsight per-query assignment with test mean loss $\le\alpha$,
unattainable since it requires knowing each answer's outcome in
advance: \system\ pays 1.4--3.2$\times$ its cost on HybridQA, CRAG, and
QAMPARI, and 6.7$\times$ on ASQA where two-constraint feasibility
forces a conservative policy --- the gap is the price of feasibility
established from finite samples rather than hindsight. The zero-loss
skyline solves a stricter objective (no loss on any query) and is shown
only for context; on three of four tracks \system\ is cheaper than it,
which says nothing about optimality --- the two solve different
problems. When the contract is infeasible for every candidate (e.g.,
$\alpha$ below the strongest plan's risk), \system\ correctly refuses to
certify and falls back to best-effort --- the planner-side analogue of
rejecting an unsatisfiable query.

\textbf{Strict contracts via certified deferral.} The headline levels
above are dictated by what always-answer pipelines can achieve; a
regulated deployment wants $\alpha \le 0.2$. Adding one knob --- defer
to human review when even the final rung's sufficiency score is below
$\kappa$ --- turns the contract into the certified pair
$\Pr(\text{answered} \wedge \text{violating}) \le \alpha$,
$\Pr(\text{deferred}) \le \beta$. At $\alpha=0.1$ \system\ certifies
HybridQA at $\beta=0.75$ (test: answered-and-violating 0.091, deferral
0.597) and ASQA at the same budget with far less deferral used (0.087,
0.219; 7.64 vs.\ 10.90 mCNY for the strongest plan); at $\alpha=0.2$
all four tracks certify --- ASQA defers only 2.9\% at $\beta=0.25$,
CRAG and QAMPARI need $\beta=0.75$ --- and at $\alpha=0.05$ no track
certifies at any tested $\beta \le 0.9$, which \system\ reports as
infeasible. Strictness is purchased transparently with deferral volume,
and the human-review load is itself certified.

\smallskip\noindent\textit{Takeaway: the guarantee costs little --- the
certified policy is 6.9--20$\times$ cheaper than the strongest plan and
within 1.4--6.7$\times$ of the matched hindsight oracle; strict
contracts down to $\alpha=0.1$ certify once deferral is priced in; every
method that beats \system\ on cost does so by breaking the contract.}

\begin{table}[t]
\caption{RQ2: single-draw test-split empirical risk / mean cost (mCNY
per query) at representative contract levels. Bold: the certified
policy, whose \emph{population} risk meets the contract with
probability $\ge 1-\delta$ (Table~\ref{tab:safety}); a single test
split carries $\pm .01$--$.04$ binomial noise and may sit on either
side of $\alpha$. The matched oracle is the cheapest hindsight
per-query assignment with test mean loss $\le\alpha$ (same objective
as \system); the zero-loss skyline solves a stricter objective and is
not cost-comparable. Full $\alpha$-grids in the artifact.}
\label{tab:main}
\centering\scriptsize
\setlength{\tabcolsep}{2.6pt}
\begin{tabular}{lcccccccc}
\toprule
& \multicolumn{2}{c}{HybridQA $\alpha{=}.34$} &
\multicolumn{2}{c}{CRAG $\alpha{=}.65$} &
\multicolumn{2}{c}{ASQA $\alpha{=}.216$} &
\multicolumn{2}{c}{QAMPARI $\alpha{=}.71$} \\
Method & risk & cost & risk & cost & risk & cost & risk & cost \\
\midrule
Fixed strongest      & .290 & 7.21 & .600 & 6.00 & .116 & 9.01 & .610 & 6.49 \\
Complexity router    & .509 & 1.46 & .620 & 4.74 & .281 & 1.14 & .612 & 6.47 \\
Utility router       & .387 & 0.40 & .625 & 0.28 & .223 & 1.02 & .712 & 0.43 \\
BO-tuned             & .380 & 0.97 & .621 & 0.89 & .221 & 1.41 & .708 & 0.94 \\
Abacus               & .321 & 1.75 & .600 & 6.00 & .201 & 1.31 & .666 & 1.23 \\
Empirical (no cert.) & .372 & 1.45 & .635 & 1.39 & .199 & 5.67 & .784 & 1.84 \\
\system\ (ladder)    & .358 & 2.11 & .618 & 2.16 & .165 & 6.42 & .724 & 2.88 \\
\system\ (full)      & \textbf{.338} & \textbf{1.05} & \textbf{.624} & \textbf{0.30} & \textbf{.145} & \textbf{6.56} & \textbf{.712} & \textbf{0.43} \\
\midrule
Matched oracle       & .340 & 0.33 & .649 & 0.18 & .214 & 0.98 & .710 & 0.30 \\
Zero-loss skyline    & .193 & 2.97 & .491 & 4.74 & .056 & 2.83 & .514 & 5.47 \\
\bottomrule
\end{tabular}
\end{table}

\begin{figure}[t]
\centering
\includegraphics[width=0.82\linewidth]{figures/cost_risk_hybridqa.pdf}
\caption{RQ2: cost--risk positions on HybridQA at $\alpha=0.34$
(test split). Certified methods (stars/triangles) sit below the contract
line; uncertified tuners that are cheaper sit above it.}
\label{fig:tradeoff}
\end{figure}

\begin{figure}[t]
\centering
\includegraphics[width=\linewidth]{figures/violation_sweep.pdf}
\caption{RQ1/RQ3: probability that the deployed policy's exact
finite-population risk exceeds $\alpha$, across 1000 calibration draws
as the contract level varies. The dotted line is the certified budget
$\delta=0.1$; only \system\ stays below it everywhere.}
\label{fig:sweep}
\end{figure}

\begin{figure}[t]
\centering
\includegraphics[width=0.7\linewidth]{figures/calibration.pdf}
\caption{RQ3: target level $\alpha$ vs.\ realized violation of the
certified policy (test split). Valid (below diagonal) but tight.}
\label{fig:calibration}
\end{figure}

\begin{figure*}[t]
\centering
\includegraphics[width=0.85\textwidth]{figures/drift.pdf}
\caption{RQ6: workload drift on CRAG ($\alpha=0.66$). Left: rolling
violation rate; the static-certified policy breaks silently after the
shift at $t{=}1500$; the e-process alarm (dash-dotted) triggers
recalibration. Right: detection delays over 20 seeded streams; no false
alarms on the 20 drift-free controls.}
\label{fig:drift}
\end{figure*}

\subsection{RQ3: calibration}
\label{sec:exp-calibration}

Figure~\ref{fig:calibration} sweeps $\alpha$ and plots realized test
violation of the certified policy. All points lie on or below the
diagonal up to test-sampling noise (maximum excess 0.009, about one
binomial standard error at $n_{\mathrm{test}}=3466$): the certificates
are valid but not vacuous --- where the contract binds, realized risk
tracks $\alpha$ with a margin of 0.02--0.07 that shrinks with calibration
size, matching the $O(\sqrt{\log(1/\delta)/n})$ radius of the HB test,
and saturates only once the cheapest plan is already feasible. The
uncertified ablation (empirical thresholds) sits \emph{on} the diagonal on
average --- and therefore above it in half the draws, which is precisely
the RQ1 failure mode.

\subsection{RQ4: rewrite vouchers}
\label{sec:exp-rewrites}

On HybridQA we fix the base plan $P^\ast$ (hybrid retrieval $k{=}16$,
rerank, qwen-plus) and apply four single rewrites, all chain prefixes
$P_1..P_4$, and the full composition on 500 calibration + 500 test
queries (Table~\ref{tab:rewrites}). A voucher is a $1-\delta_r$
confidence bound, so the correct empirical test is \emph{coverage}:
pooling cal+test into a fixed 1000-query population and redrawing the
500-query calibration set 1000 times, each single-rewrite voucher
(Clopper--Pearson, $\delta_r{=}.05$) covers the population's exact harm
rate in 98.9--99.5\% of draws, above its 95\% nominal level --- the
one-sided bound's expected conservatism. (A single test-split column
can land on either side of the voucher, as R1 and R2 do in the table;
that is binomial noise, not a coverage failure.) The harm ordering
(prefilter $<$ rerank-elision $\approx$ $k$-truncation $<$ generator
downgrade) matches the intuition that evidence-preserving rewrites are
safest. The composed plan cuts cost 8.3$\times$ (1.74 $\to$ 0.21 mCNY)
with population harm 0.240.

\textbf{Composition.} The base-relative union bound (0.342) requires the
no-synergistic-harm assumption --- and the data \emph{violate} it: 8\% of
harmed queries are harmed only by the composition (synergy between
truncation and generator downgrade), so that bound's nominal coverage is
unfounded even where its number happens to be large enough. The
telescoping path voucher of Proposition~\ref{prop:composition}
(steps $.026/.081/.105/.164$, sum 0.376) costs only 10\% extra width
over the union bound while being assumption-free --- covering the
population chain harm in every one of the 1000 redraws ($>$99.7\% at
95\% confidence) --- and both remain
looser than the direct certificate (0.251 mean, 98.9\% coverage) on the
deployed chain --- which is why \system\ uses vouchers to order search
and certifies directly.

\textbf{Selectivity-conditioned vouchers.} The flat prefilter voucher
($\hat\eps = .026$) is marginally valid but hides heterogeneity that the
planner's selectivity estimate exposes: stratifying by pool selectivity
$\sigma$, realized test harm is 0.075 on the vague-predicate stratum
($\sigma \in [.79, .91)$) --- 3$\times$ the flat voucher --- versus
none observed on the $\sigma = 1$ stratum. The per-stratum vouchers
($.076/.088/.128/.020$, Bonferroni $\delta_r/4$) cover every stratum's
realized harm, giving the rewrite search a risk signal that tracks
\emph{which queries} a pushdown endangers rather than an average.

\begin{table}[t]
\caption{RQ4: rewrite vouchers on HybridQA ($n_{\mathrm{cal}}{=}500$;
cost in mCNY per query).
Harm = loss strictly worse than the base plan's. Vouchers $\hat\eps$
are $1-\delta_r$ confidence bounds; ``cover.'' = fraction of 1000
calibration redraws whose voucher covers the exact harm rate of the
pooled 1000-query population (nominal 95\%). The three composition
rows price the \emph{same} composed plan, so their execution columns
coincide by construction and only the bound differs; the telescoping
voucher covered in every one of the 1000 redraws, reported as the
one-sided 95\% lower bound $>$.997 against its nominal $.80$ level
($m{=}4$ steps at $\delta_r{=}.05$) --- the deliberate price of
dropping the NSH assumption, which the workload violates
(8\% synergistic harm); the telescoping path voucher
(Prop.~\ref{prop:composition}) is assumption-free at 10\% extra width
over the (invalid) union bound.}
\label{tab:rewrites}
\centering\scriptsize
\setlength{\tabcolsep}{2.2pt}
\begin{tabular}{lcccccc}
\toprule
Variant & cost & harm$_{\mathrm{cal}}$ & voucher $\hat\eps$
& harm$_{\mathrm{te}}$ & harm$_{\mathrm{pop}}$ & cover. \\
\midrule
base $P^\ast$    & 1.74 & ---  & ---  & ---  & --- & --- \\
R1 prefilter     & 1.75 & .014 & .026 & .030 & .022 & .995 \\
R2 $k$-trunc     & 1.11 & .060 & .081 & .084 & .072 & .991 \\
R3 no-rerank     & 1.69 & .058 & .078 & .070 & .064 & .994 \\
R4 gen-downgr.   & 0.33 & .130 & .157 & .152 & .141 & .989 \\
\midrule
comp.\ (union, NSH) & 0.21 & .218 & .342 & .262 & .240 & --- \\
comp.\ (telescoping) & 0.21 & .218 & .376 & .262 & .240 & $>$.997 \\
comp.\ (direct cert.) & 0.21 & .218 & .251 & .262 & .240 & .989 \\
\bottomrule
\end{tabular}
\end{table}

\subsection{RQ5: group-conditional contracts}
\label{sec:exp-group}

CRAG's dynamism labels expose what marginal guarantees hide. With
group-wise contract levels
$\alpha_{\text{static}} = .62$, $\alpha_{\text{slow}} = .72$,
$\alpha_{\text{fast}} = .92$, and
$\alpha_{\text{rt}} = .93$, a single marginally-certified
policy (at the mixture level) violates the real-time group's contract in
\textbf{99\%} of 200 draws (realized risk 0.99 vs.\ 0.93 target: the cheap
rungs it certifies for the average workload have no fresh evidence for
real-time queries) and the static group's in 14.5\%. Group-conditional
certification (Bonferroni $\delta/4$ per group) brings every group to
$\le 3.5\%$, escalating real-time queries to the KG-API rung (per-query
cost 7.3 mCNY on that 9\% slice) while \emph{lowering} static-group cost.
When a group's contract is infeasible in the plan space --- fast-changing
at $\alpha=.65$, where even the strongest plan's group risk is 0.79 ---
\system\ reports non-certifiability instead of silently under-delivering.
For Example~\ref{ex:finance} this is the difference between a compliance
report that says ``real-time market questions cannot currently meet the
freshness bar'' and one that quietly averages the failure away.

\subsection{RQ6: drift detection and recovery}
\label{sec:exp-drift}

\textbf{Workload drift.} A stream of 1500 static/slow queries followed by
1500 fast/real-time queries (real losses, designed arrival order) breaks
the marginally-certified policy: violation rises from 0.58 to 0.88
(contract $\alpha=0.66$), and without monitoring it never knows
(Fig.~\ref{fig:drift}). The e-process monitor fires at a median of 78
queries after onset (range 36--125 over 20 seeds) with \textbf{zero} false
alarms across 20 drift-free control streams --- an expected zero, not a
suspicious one: the per-stream Ville bound is $\delta = 0.1$, and on a
compliant stream the supermartingale in fact decays, so the realized
false-alarm rate sits far below the bound (confirmed at 0/200 in the
magnitude study below).
Post-alarm, \system\ switches to group-conditional certified policies
(strongest-plan best effort plus an infeasibility flag for the fast-changing
group), cutting phase-2 violation to 0.82 --- the residual gap is the
infeasible slice, now \emph{reported} rather than hidden.

\begin{figure}[t]
\centering
\includegraphics[width=\columnwidth]{figures/drift2.pdf}
\caption{RQ6: measured detection delay (mean and p90 over 30 seeds per
drift magnitude) versus the Proposition~\ref{prop:delay} rate
$(\log(1/\delta) + \log(T|W|))/g^\ast$ --- mean delay sits below the
rate at every magnitude and the p90 tracks it within 0.6--1.0$\times$
(left). The $O(1)$ CUSUM
recursion localizes the changepoint to within a few queries for
strong drifts (right). No false alarms on the 200 drift-free
controls.}
\label{fig:drift2}
\end{figure}

\textbf{Delay scaling and changepoint localization.} Varying the drift
magnitude by mixing pools (post-change dynamic fraction
$\rho \in \{.5, .75, 1\}$, 30 seeds each) makes
Proposition~\ref{prop:delay} quantitatively checkable: for each
magnitude we compute the post-change optimal log-growth $g^\ast$ on the
monitor's bet grid ($.0071/.049/.137$) and compare measured delays
($760/203/82$ queries mean) against the rate
$(\log(1/\delta) + \log(T|W|))/g^\ast = 15.4/g^\ast$. Mean delay sits
below the rate at every magnitude, is unaffected by the 1500-query
compliant prefix, and the rate is \emph{tight}, not ornamental: the
p90 delay reaches 0.91--1.04$\times$ of it at moderate and strong
drifts (Fig.~\ref{fig:drift2}) --- the analysis predicts the
detector's operating curve, not just its order of magnitude. False
alarms: none across the 200
drift-free control streams ($<1.5\%$ at 95\% confidence; Ville's
inequality is
conservative in practice). The $O(1)$ CUSUM recursion localizes the
changepoint to a median error of \emph{4 queries} at $\rho{=}1$ and 10
queries at $\rho{=}.75$ --- sharp enough that recalibration can discard
exactly the pre-change prefix.

\textbf{Closed-loop certified recalibration.} An alarm is not yet a
repaired policy. Over 10 seeded streams we run the full loop: alarm
$\to$ safe mode (strongest rung --- a best-effort fallback that carries
\emph{no} certificate on the unseen post-change distribution, and is
flagged as such) for a 250-query labeled window $\to$ fixed-sequence
recertification on that window $\to$ deploy the newly certified policy,
which is where the post-drift guarantee resides. The loop holds phase-2
violation to
0.819 versus 0.886 for the static policy, within half a point of the
oracle that switches to the post-change-certified policy at the exact
drift time (0.815), and beats the DB-style periodic reoptimizer
(recheck every 250 queries: 0.837) --- while spending less than the
oracle (4.77 vs.\ 4.94 mCNY/query) because the certified post-alarm
policy backs off from the safe mode as soon as the window certifies a
cheaper ladder.

\textbf{Lifetime operation across repeated alarms.}
Proposition~\ref{prop:lifetime} prices an unbounded alarm--recalibrate
loop by geometric budget spending ($\delta_k = \delta/2^{k+1}$ per
epoch); we exercise it on a four-phase stream (static $\to$ real-time
$\to$ static $\to$ fast, 1200 queries each, $\alpha=0.85$, 20 seeds).
Runs average 1.75 alarms; the first drift is caught in a median of 98
queries and the second --- three epochs later, on a quarter of the
initial budget --- in 801, the logarithmic delay price the proposition
predicts. Recalibration certifies a post-change policy in 84\% of alarm
epochs; in the remainder (the real-time phase, where even the strongest
plan is out of contract) the system reports non-certifiability and
stays in flagged safe mode. Scheduled quiet-period rechecks recover
cheap policies after drifts revert (post-drift static phase: violation
0.574 vs.\ 0.603 up to complete failure --- every real-time-phase
query out of contract --- for never recalibrating). No false alarms
occurred
across 100 null streams ($<3\%$ at 95\% confidence) despite five
sequential tests per run.

\textbf{Audit accounting and subsampled monitoring.} The monitor
consumes audited violation indicators, which are not free: on CRAG,
exact-match short answers resolve 44\% of audits at negligible cost and
the rest need a qwen-max judge call at 0.52 mCNY --- an effective 0.29
mCNY per audited query, comparable to the certified policy's own 0.30
mCNY deployment cost. Charged honestly, at 100\% audit the certified
system's total cost of ownership (deployment + audit + amortized
calibration) is 0.60 mCNY/query --- RQ2's ``20$\times$'' saving is
10.0$\times$ with the monitor metered in. Auditing an independent
random $r$-fraction of traffic keeps every guarantee on the audited
subsequence (the audited stream is an unbiased draw;
\S\ref{sec:monitor}) and dilates delay by $1/r$: measured medians
74/343/878 queries at $r = 1/0.25/0.1$ with no false alarms at any
rate, restoring savings to 15.8$\times$ and 17.9$\times$. The audit
rate is thus an explicit dial between reaction time and monitoring
spend.

\begin{figure*}[t]
\centering
\includegraphics[width=0.85\textwidth]{figures/stale.pdf}
\caption{RQ6b: index staleness on CRAG (joint correct-and-fresh contract,
$\alpha=0.75$). Left: the index silently falls one week behind at
$t{=}2000$; every plan's joint risk exceeds $\alpha$, the static-certified
policy breaks, the e-process alarm triggers a switch to the
stale-state-certified policy (freshness-honest re-execution, KG-API
escalation). Right: detection delays over 20 seeded streams; no false
alarms on the 20 in-sync controls.}
\label{fig:stale}
\end{figure*}

\textbf{Index staleness.} We re-execute the CRAG ladder (real retrieval
and generation; expanded cal/test of 1000/1006 queries per index state)
under a 72h freshness contract in two index states: in-sync, and one week
behind. The joint loss is
$\mathbf{1}[\text{not correct} \vee \ell_{\text{fresh}}]$. When the index
falls behind \emph{silently}, the executor keeps serving dated web
evidence that is actually 168h older than it believes: per-rung joint risk
jumps from $(.79, .74, .70, .70)$ to $(.94, .91, .91, .90)$ --- even the
strongest plan now violates $\alpha=0.75$, so no amount of static
escalation helps (Fig.~\ref{fig:stale}). The static-certified policy runs
at phase-2 violation 0.926 and never knows; the monitor alarms after a
median of 117 queries
(range 38--240 over 20 seeds, no false alarms on the 20 in-sync
controls)
and switches to the policy certified on stale-state calibration data ---
a freshness-honest re-execution whose filter knows the lag, drops
out-of-contract evidence, and escalates to the still-fresh KG API ---
restoring violation to 0.759~$\approx\alpha$ at 2.52 mCNY/query, within
1\% of the oracle that recalibrates exactly at the drift point (0.752 at
2.55). Freshness contracts are thus operationally different from quality
contracts: staleness breaks \emph{every} plan at once, and only
detect-and-recertify --- not conservative planning --- restores the
guarantee.

\subsection{RQ7: multi-constraint contracts and optimizer ablation}
\label{sec:exp-multi}

Joint contracts expose the value of a heterogeneous certified policy
space. On ASQA with citation $\alpha_q{=}.3$ \emph{and} a 10s latency SLO
($\alpha_\ell{=}.15$): quality-only certification meets citations but
violates the SLO on 24.8\% of queries (escalation accumulates latency);
the escalation-only family has \emph{no} feasible candidate; the full
optimizer certifies a jump router that satisfies both constraints
(realized citation risk 0.230; no SLO violation observed on test) at
0.98 mCNY --- 5.0$\times$ cheaper than the
quality-only policy --- because jumping pays a single rung's latency.

\textbf{Four-constraint contracts.} The full contract of
Example~\ref{ex:finance} binds four dimensions at once. On CRAG we
certify correctness $\le .65$, hallucination $\le .30$, staleness
$\le .05$, and latency-budget exceedance $\le .10$ (8s) jointly: the
IUT-certified policy satisfies all four on test (correctness .603,
hallucination .258, latency .057, every per-constraint p-value
$\le 5.5{\times}10^{-3}$; the staleness loss never fires on the
1006-query test split, $<.003$ at 95\% confidence,
because the certified policy escalates every stale-evidence query) at
5.02 mCNY/query, certifying 47 candidates
against Bonferroni's 44 and selecting a 3\% cheaper policy. On ASQA
(citation recall $\le .25$, citation precision $\le .35$, quality
$\le .40$, latency $\le .15$ at 15s) the gap is qualitative: IUT
certifies 45 candidates (test risks $.188/.230/.261/.033$) while
Bonferroni certifies \emph{none} --- at $n_{\mathrm{cal}}{=}350$ the
$\delta/4$ split is fatal, exactly as \S\ref{sec:certification}
predicts.

\textbf{Multi-constraint drift attribution.} Post-deployment, four
e-process monitors (one per constraint, each at $\delta_{\mathrm{mon}}/4$,
so the familywise false-alarm budget is one global
$\delta_{\mathrm{mon}}$) watch the certified CRAG policy. Under a
fresh-to-stale index transition the freshness monitor alarms after a
median of 13.5 queries (range 5--25, 20 seeds) while the correctness,
hallucination, and latency monitors --- whose constraints genuinely
remain in contract --- stay silent; all four are silent on 40
drift-free controls. The alarm is thus also an \emph{attribution}: the
operator learns that the index, not the generator, is failing.

\smallskip\noindent\textit{Takeaway: multi-dimensional contracts of the
kind in Example~\ref{ex:finance} are exactly where the heterogeneous
policy space and the IUT earn their keep --- Bonferroni certifies
nothing on ASQA where the IUT certifies 45 candidates, and per-constraint
monitors under one global budget turn alarms into attributions.}

\subsection{RQ8: generalization across model families and serving
backends}
\label{sec:exp-families}

\begin{figure*}[t]
\centering
\includegraphics[width=0.95\textwidth]{figures/families.pdf}
\caption{RQ8: the same certified optimizer across four generator
families (HybridQA, $\alpha{=}0.35$; gemma at $\alpha{=}0.45$).
Retrieval is materialized once and shared; only the generator column
changes. Grey: the four fixed rungs; star: the certified policy.
The optimizer picks a \emph{different} policy family per price
structure (threshold ladder for Qwen, jump routers for DeepSeek/GLM/
Gemma) and stays below the contract line everywhere it certifies.}
\label{fig:families}
\end{figure*}

\begin{table*}[t]
\caption{RQ8: the same certified optimizer across four generator
families from four model developers, served through two backends ---
Qwen, DeepSeek, and GLM hosted on one API platform (DashScope,
token-metered at each developer's list prices), Gemma-3 served locally
by Ollama and metered in GPU-seconds (HybridQA, quality contract
$\mathbf{1}[\text{F1}<0.5]$, $\delta{=}0.1$; splits 300/500/600 over
shared materialized retrieval). ``Viol.\ prob.'' = fraction of 200
repeated calibration/test draws whose exact finite-population risk
exceeds $\alpha$ (must be $\le \delta$); ``emp.'' = the
empirical-no-certificate ablation; cells with no violating draw report
the one-sided 95\% upper bound $<.015$ --- the protocol's resolution
at 200 draws --- the expected outcome under the
certificate's built-in margin; held-out rates are 4.5\% (Qwen) and
9.5\% (DeepSeek). \system\ is the only method within
budget on every family; the certified policy family itself adapts to
each developer's price structure. At $\alpha{=}0.35$ the Gemma-3 plan
space is infeasible and \system\ reports non-certifiability instead of
deploying the least-bad plan.}
\label{tab:families}
\centering\footnotesize
\setlength{\tabcolsep}{3pt}
\begin{tabular}{llcccccccccc}
\toprule
& & & \multicolumn{2}{c}{Strongest fixed plan}
& \multicolumn{4}{c}{\system\ (certified)}
& \multicolumn{2}{c}{Utility router} & Oracle \\
\cmidrule(lr){4-5}\cmidrule(lr){6-9}\cmidrule(lr){10-11}
Family (serving) & $\alpha$ & metering
& risk & mCNY & risk & mCNY & saving & selected policy
& risk & mCNY & mCNY \\
\midrule
Qwen (DashScope)     & .35 & tokens & .300 & 7.29 & \textbf{.320} & 6.94 & 1.1$\times$ & ladder $q{=}.85$ & .395 & 0.37 & 3.18 \\
DeepSeek (DashScope) & .35 & tokens & .295 & 33.51 & \textbf{.337} & 2.02 & \textbf{16.6$\times$} & router $\lambda{=}32$ & .365 & 1.32 & 12.84 \\
GLM (DashScope)      & .35 & tokens & .277 & 22.71 & \textbf{.317} & 4.06 & 5.6$\times$ & router $\lambda{=}100$ & .355 & 3.76 & 13.68 \\
Gemma-3 (local Ollama) & .45 & GPU-s & .395 & 0.66 & \textbf{.432} & 0.57 & 1.2$\times$ & router $\lambda{=}178$ & .465 & 0.49 & 1.06 \\
\midrule
\multicolumn{3}{l}{Viol.\ prob.\ over 200 draws (cert.\ / emp.)}
& \multicolumn{9}{l}{Qwen \textbf{$<$.015}/.300 \quad DeepSeek \textbf{$<$.015}/.095 \quad GLM \textbf{$<$.015}/.055 \quad Gemma-3 \textbf{$<$.015}/$<$.015 \qquad (population risk; $\delta{=}0.1$)} \\
\bottomrule
\end{tabular}
\end{table*}

A fair objection to any cost-optimization result is that it may be an
artifact of one vendor's price sheet. We therefore hold the entire
optimizer fixed --- retrieval ladder, sufficiency scorers, candidate
families, certification --- and swap only the generator column across
four model families ---
Qwen~\cite{qwen25}, DeepSeek~\cite{deepseekv3,deepseekr1},
GLM~\cite{glm2024}, and open-weight Gemma-3~\cite{gemma3} ---
(Table~\ref{tab:families},
Fig.~\ref{fig:families}; splits 300/500/600
train/cal/test over shared materialized retrieval). What varies is the
model developer and its price structure (four developers, two serving
backends: one hosted API platform plus one local deployment); testing
robustness across independent hosting platforms for the \emph{same}
model would require replicating the ladder on a second API vendor,
which we leave outside scope and do not claim. Four results.
\emph{(i) The guarantee transfers}: on every family where the contract
is feasible, the certified policy's realized test risk is below
$\alpha$ --- Qwen 0.320, DeepSeek 0.337, GLM 0.317 at $\alpha{=}0.35$;
Gemma-3 0.432 at $\alpha{=}0.45$ --- and over 200 repeated
calibration/test draws per family (RQ1's population protocol on the
pooled 1{,}100 queries) the certified policy's exact population risk
never exceeds $\alpha$ in \emph{any} draw (below 1.5\% at 95\%
confidence on all four families
against the $\delta{=}0.1$ budget; up to 9.5\% on the noisier held-out
criterion), while the
empirical-no-certificate ablation deploys truly out-of-contract
policies in up to \textbf{30\%} of
draws (Qwen; 9.5\% DeepSeek, 5.5\% GLM) and the tuned utility router's
single deployment exceeds $\alpha$ on \emph{all four} families
(risk .355--.465). Uncertified selection does not fail on one vendor's
quirk; it fails everywhere.
\emph{(ii) The savings are not
Qwen-specific --- they grow with the family's price spread}: against
each family's strongest plan, certified cost drops 16.6$\times$ for
DeepSeek (2.0 vs.\ 33.5 mCNY/query), 5.6$\times$ for GLM (4.1 vs.\
22.7), 1.1--1.2$\times$ for Qwen and Gemma at these $\alpha$ (the Qwen
figure is smaller than RQ2's 6.9$\times$ because $n_{\mathrm{cal}}$
here is 500, not 3000: the wider certificate radius forces a deeper
stopping rule --- the price of a smaller audit budget, not a property
of the family); and the
\emph{certified policy family itself changes} with the price structure
--- a threshold ladder for Qwen's smooth tiering, jump routers for
DeepSeek and GLM, whose flash tiers are strong but whose top tiers are
disproportionately expensive. The optimizer discovers this from data;
nothing in the system encodes a vendor. \emph{(iii) Infeasibility is
reported, not hidden}: at $\alpha{=}0.35$ no Gemma-3 plan --- not even
the 27b rung (risk 0.395) --- can meet the contract, and \system\
returns non-certifiable instead of deploying the least-bad plan, the
same honest failure mode as RQ5. \emph{(iv) Cost-metric generality}:
Gemma plans are metered in measured GPU-seconds (0.66 mCNY/query at
the 27b rung --- an order of magnitude below hosted tiers at 10 points
more risk), and certification is indifferent to which currency the
cost column uses. Re-certifying over the 8{,}029-candidate Cartesian
grid instead of the 69-candidate default changes validity nowhere
(selection-proofness) but illustrates a real power effect: at
$n_{\mathrm{cal}}{=}500$ the grid's cross-fit walk order reaches the
contract boundary sooner on Qwen and certifies nothing, falling back
to the strongest plan --- more candidates are statistically free for
\emph{safety} but not automatically better for \emph{power} at small
audit budgets, which is why \system\ ships the compact curated space
as the default.

\subsection{RQ9: access-path selection and plan-space scale}
\label{sec:exp-system}

\textbf{Join ordering / access paths.} On HybridQA we execute four
physical plans for the table$\bowtie$text join at a \emph{fixed}
generator (qwen-flash), so plan differences are purely retrieval-side:
\{table-first, text-first\} $\times$ \{lexical, hybrid+rerank\}, with
GPU-priced retrieval included in cost. Estimated pool selectivity
$\sigma$ (computed from the row-link cardinalities \emph{before}
execution, \S\ref{sec:planner}) stratifies violation cleanly and
monotonically for every path: e.g., text-first-hybrid moves from 0.511
(lowest-$\sigma$ tercile) to 0.364 (highest), and the lexical paths pay
a violation premium over hybrid that \emph{grows} with $\sigma$ (1.9
points on the vague-predicate tercile, where even reranking cannot
recover a poor pool, versus 5.1 on the selective one) --- the
selectivity-dependent crossover that access-path selection exploits.
Feeding the chooser family (fixed paths + selectivity-threshold rules
over a target grid) to the certification layer at $\alpha{=}0.5$
certifies a mixed policy at 0.337 mCNY/query --- 31\% below the
strongest fixed path (text-first-hybrid, 0.492) with realized test
violation 0.481 $\le \alpha$; the per-query oracle path (0.374 at
0.347) bounds the headroom better cardinality estimates could still
buy. Everything a relational optimizer does with selectivities ---
pick the join order, skip the expensive operator when the predicate is
strong --- transfers, except that feasibility is established by the
test, not by the estimate.

\textbf{Plan-space scalability.} Certification is a vectorized pass
over the materialized ladder matrix, so enlarging the plan space is
cheap: growing the threshold grid from 69 to 97{,}365 candidates
(Cartesian per-rung thresholds, the $10^5$ regime of combinatorial
plan enumeration) raises certification wall-clock to 8.5s on HybridQA,
4.1s on CRAG, and 0.9s on ASQA --- still below the cost of a single
qwen-max call, with selection-proofness (Theorem~\ref{thm:validity})
guaranteeing that the extra $10^3\times$ candidates cost nothing in
validity. \textbf{Operator sharing.} Materializing retrieval once and
sharing it across the four generator families of RQ8 removes
three-quarters of retrieval compute (8--19\% of single-family
wall-clock, a measured 1.06--1.17$\times$ end-to-end speedup for
$F{=}4$) and, more importantly, eliminates GPU contention between the
reranker and locally served generators.
